\section{Conclusion}
\label{sec:conclusion}

Subject awareness in EEG denoising does not require learning the subject into the network. The quantity that actually differs across people, the propagation of ocular potentials to the scalp, can be measured from two minutes of paired EEG and EOG --- one minute, our duration curve shows, already buys 97\% of the benefit --- stabilized by shrinkage, checked for reliability, and handed to a population diffusion model as an explicit artifact estimate. On this design the individual estimate improved reconstruction for every development participant and reproduced at magnitude in a single-pass held-out evaluation; on natural ambulatory recordings it removed more ocular signal while preserving more of the recording than population calibration, and it is the only method in our comparison, classical anchors included, on the healthy side of the attenuation--retention plane; and the same calibration uncertainty, propagated into the sampler, produced predictive intervals with nominal 80\% and 90\% coverage that held at 0.81 and 0.86 on held-out users under a temperature frozen in advance. Downstream, the cleaned signal decodes as well as the raw recording, and where unguided cleaning costs task information the calibrated guide restores it. The failure modes are measured parts of the system rather than caveats around it: a rejected calibration withholds the guide, because substituting a population estimate subtracts the wrong signal; a wrong person's calibration is harmful, so the protocol keeps calibration and use in the same session; the calibrated operator has a measured lifetime, so adaptation runs through recalibration --- and on-line tracking during use never reaches the calibrated operating point.

The evaluated scope is same-session ocular-artifact removal with recorded EOG and a compatible montage. Natural recordings supply attenuation and preservation measures, not a clean target; the held-out cohort has eight participants, mitigated by the frozen single-pass protocol rather than by sample size, and a sealed 55-participant block on a second dataset is reserved for a future confirmatory evaluation; as a point estimator the calibrated diffusion model matches calibrated linear regression, and its distinctive contributions are the predictive distribution, the robustness of one network across all operating modes, and the measured adaptation loop around them. The supporting experiment suggests the design generalizes past its own interface: calibrated subject conditioning also works through a learned representation, flat in training-population size --- so the road ahead for subject-aware diffusion denoising is not larger population models but better calibration channels. External EEG--EOG cohorts, non-ocular artifact families, and compute-matched distributional baselines are the natural extensions. Within the evaluated setting the conclusion is direct: the operator carries the subject, and calibrating it adapts a shared generative denoiser --- accurately, safely, and with honest uncertainty --- to a person the model has never seen.
